Krátký kurátorský přehled zdrojů, tutoriálů a doporučených kroků pro další vzdělávání pedagogů při pilotním nasazení generativní AI ve výuce. Níže najdete ověřené praktické vstupy (nástroje a ukázky aktivit) s odkazy na doporučené postupy a krátký akční plán „co dělat dál“.

Doporučené online tutoriály a vstupní zdroje
\begin{itemize}
  \item Minecraft Education — úvodní interaktivní světy a Hour of Code jako jednoduchý hands‑on způsob, jak studentům přiblížit principy strojového učení a agentního chování \cite{kb_ai_101_tipu_pdf_04e4a115_page_15_2}. (Praktické: stáhnout stránky Minecraft Education a vyzkoušet jednu z volně dostupných hodin.)
  \item Pokrok v matematice (Microsoft educational accelerators) — demonstrace generování sady příkladů, sběru postupů a předběžného opravení studentských odpovědí nástrojem; užitečné pro cvičení a diagnostiku slabých míst ve třídě \cite{kb_ai_101_tipu_pdf_04e4a115_page_8_1}.
  \item Microsoft 365 Copilot — pro úlohy „deep reasoning“ doporučen agent Výzkumník (Researcher), který v rámci Copilota zpracovává rozsáhlejší souvztažnosti a podklady; vhodné při přípravě strategických dokumentů a komplexnějších výukových materiálů \cite{kb_ai_101_tipu_pdf_04e4a115_page_48_1}.
  \item Repozitář VUT s šablonami a video‑tutoriály — živý repozitář (ai.vut.cz) obsahující šablony promptů, checklisty pro piloty a krátké video‑návody; doporučený další krok po ukončení lokálního pilotu (kontaktujte metodické oddělení pro přístup a podporu). (Organizační informace vyplývají z interních doporučení příručky — viz úvodní část příručky.)
\end{itemize}

Rychlý akční plán pro pedagoga — 5 kroků k prvnímu pilotu
\begin{enumerate}
  \item Projděte si krátké video‑tutoriály v repozitáři ai.vut.cz a vyberte jednu nízkorizikovou aktivitu (např. Minecraft Education nebo Pokrok v matematice). (repozitář jako další krok po pilotu; viz výše)
  \item Před spuštěním aktivity vypracujte minimální DPIA / datový plán a ověřte, že se do služby neodesílají osobně identifikovatelné údaje; zajistěte pseudonymizaci, pokud je to nutné \cite{kb_malinka_chatgpt_ve_vyuce_dva_roky_pote_2024_pdf_90bc3aaf_page_12_1}.
  \item Spusťte krátký pilot (30–90 min) s jasně definovanými cíli a metrikami (pedagogický cíl, očekávaný výstup, způsob validace). Jako první piloty doporučujeme aktivity popsané výše: Minecraft Education (Hour of Code) a Pokrok v matematice \cite{kb_ai_101_tipu_pdf_04e4a115_page_15_2,kb_ai_101_tipu_pdf_04e4a115_page_8_1}.
  \item Dokumentujte artefakty: prompty, varianty odpovědí, záznamy o opravách a krátkou reflexi (co fungovalo, kde AI selhala). Ponechte vždy human‑in‑the‑loop — vyučující kontroluje a finalizuje výstupy \cite{kb_malinka_chatgpt_ve_vyuce_dva_roky_pote_2024_pdf_90bc3aaf_page_12_1}.
  \item Pokud projekt vyžaduje hlubší rešerši nebo strategickou přípravu obsahu, vyzkoušejte použití Microsoft 365 Copilot s agentem Výzkumník pro agregaci a strukturování podkladů \cite{kb_ai_101_tipu_pdf_04e4a115_page_48_1}.
\end{enumerate}

Krátký checklist před pilotem (rychlé bezpečnostní body)
\begin{itemize}
  \item Existuje minimální DPIA / datový plán pro aktivitu? (ano / ne) \cite{kb_malinka_chatgpt_ve_vyuce_dva_roky_pote_2024_pdf_90bc3aaf_page_12_1}
  \item Jsou vyloučena osobně identifikovatelná data z odesílaného obsahu? (ano / ne) \cite{kb_malinka_chatgpt_ve_vyuce_dva_roky_pote_2024_pdf_90bc3aaf_page_12_1}
  \item Je vymezeno, kdo finalizuje výstupy a zajišťuje validaci (human‑in‑the‑loop)? (ano / ne) \cite{kb_malinka_chatgpt_ve_vyuce_dva_roky_pote_2024_pdf_90bc3aaf_page_12_1}
  \item Máte zdokumentované prompty a krátký plán pro evaluaci pedagogického přínosu? (ano / ne)
\end{itemize}

Kde najít další materiály
\begin{itemize}
  \item VUT repozitář (ai.vut.cz) — šablony promptů, checklisty pro DPIA a krátké video‑návody (living document; kvartální aktualizace).
  \item Ověřené ukázky aktivit v dokumentaci Microsoft Education (Minecraft Education, Pokrok v matematice) pro rychlé hands‑on nasazení \cite{kb_ai_101_tipu_pdf_04e4a115_page_15_2,kb_ai_101_tipu_pdf_04e4a115_page_8_1}.
  \item Pro hlubší rešerši a tvorbu strategických materiálů doporučujeme prozkoumat funkce agenta Výzkumník v Microsoft 365 Copilot \cite{kb_ai_101_tipu_pdf_04e4a115_page_48_1}.
\end{itemize}

Poznámka (general background knowledge): krátké tutoriály a praktická cvičení s jasně vymezenými bezpečnostními pravidly jsou nejrychlejší cestou, jak si vybudovat zkušenosti a získat podklady pro širší zavedení AI na úrovni katedry.